\documentclass{article}
\usepackage{graphicx} % Required for inserting images

% Added stuff
\usepackage{amsmath}
\usepackage{amssymb}
\newcommand{\N}{\mathbb{N}}
\newcommand{\Z}{\mathbb{Z}}
\newcommand{\Q}{\mathbb{Q}}
\newcommand{\R}{\mathbb{R}}
\newcommand{\C}{\mathbb{C}}
\newcommand{\bigp}[1]{\left(#1\right)}
\newcommand{\bigb}[1]{\left[#1\right]}
\newcommand{\bigabs}[1]{\left|#1\right|}
\newcommand{\summation}[1][0]{\sum_{n=#1}^\infty}
\newcommand{\summationk}[1][0]{\sum_{k=#1}^\infty}
\newcommand{\summationm}[1][0]{\sum_{m=#1}^\infty}

\title{Infinite Series 5}
\author{Siddharth K}
\date{February 2026}

\begin{document}

\maketitle

\section*{Problem 1}
Take any $\alpha = a + b\sqrt {-d}$ with $a, b \in \Q$. For $\alpha$ to be integral, $(x - \alpha)(x - \bar\alpha)$ must have integer coefficients. We have that
$$(x - \alpha)(x + \alpha) = (x - (a + b\sqrt{-d}))(x - (a - b\sqrt{-d}))$$
$$= x^2 - 2ax + a^2 + db^2,$$
so $2a$ and $a^2 + db^2$ must be integers. Let $u = 2a$ and $v = 2b$. We have that $u$ must be an integer. Now, consider $4(a^2 + db^2)$:
$$4(a^2 + db^2) = (2a)^2 + d(2b)^2 = u^2 + dv^2.$$
That must be an integer, and since $u$ is an integer, $v = 2b$ must be an integer as well. So every integral element has the form
$$\alpha = \frac{u + v\sqrt{-d}}{2},\quad u, v \in \Z.$$
Now, we need to find which $u, v$ give integers. For $a^2 + db^2$ to be an integer,
$$\frac{u^2 + dv^2}{4} \in \Z \iff u^2 \equiv -dv^2 \pmod{4}.$$
If $v$ is even, then $v^2 \equiv 0 \pmod{4}$, so $u$ must be even. If $v$ is odd, then $v^2 \equiv 1 \pmod{4}$, so $u^2$ is not even and $u^2 \equiv 1 \pmod{4}$. So we need $-d \equiv 1 \pmod{4}$. If $-d \not\equiv 1 \pmod{4}$, no integral element can have $v$ odd, so $\alpha$ must lie in $\Z[\sqrt{-d}]$.

\section*{Problem 5}
Define multiplication on classes by $[a] \cdot [b] = [ab]$. \\
We need to show that if $a \sim a'$ and $b \sim b'$ then $ab \sim a'b'$. Assume that $(c)a = (d)a'$ and $(e)b = (f)b'$ for some $c, d, e, f \in R$. Then
$$(ce)(ab) = ((c)a)((e)b) = ((d)a')((f)b') = (df)(a'b').$$
So $(ce)(ab) = (df)(a'b')$, so $ab \sim a'b'$.
Next, we need to show associativity. We know that $(ab)c = a(bc)$, so
$$([a] \cdot [b]) \cdot [c] = [ab] \cdot [c] = [(ab)c] = [a(bc)] = [a] \cdot [bc] = [a] \cdot ([b] \cdot [c]).$$

\section*{Problem 8}
Each term in the product is the sum of a geometric series:
$$\prod_{p \text{ prime}} \frac{1}{1 - \frac{\chi(p)}{p^s}} = \prod_{p \text{ prime}} \summation \bigp{\frac{\chi(p)}{p^s}}^n.$$
Expanding the product term-by-term, it sums over each unique prime factorization, which is just all natural numbers, and $\chi$ is multiplicative, so the product is
$$\summation[1] \frac{\chi(n)}{n^s}.$$

\end{document}